\documentclass[11pt,a4paper]{article}

\usepackage[utf8]{inputenc}
\usepackage[T1]{fontenc}
\usepackage[french]{babel}
\usepackage[margin=2.2cm]{geometry}
\usepackage{booktabs}
\usepackage{tabularx}
\usepackage{array}
\usepackage{float}
\usepackage{enumitem}
\usepackage[table]{xcolor}
\usepackage[colorlinks=true,linkcolor=black,urlcolor=blue]{hyperref}

\newcolumntype{Y}{>{\raggedright\arraybackslash}X}
\setlength{\parindent}{0pt}
\setlength{\parskip}{0.5em}
\renewcommand{\arraystretch}{1.25}

\title{Feuille de route de développement\\[0.3em]\large Projet C++ \& Software Engineering --- Éditeur de maillage 3D (Half-Edge)}
\author{Ibrahima Sory Diakite}
\date{Année universitaire 2025--2026}

\begin{document}
\maketitle

\section{Objectif et principes directeurs}

Le projet consiste à concevoir un éditeur de maillage 3D interactif en C++ moderne, fondé sur la structure Half-Edge canonique à quatre pointeurs par demi-arête (\texttt{vertex}, \texttt{face}, \texttt{next}, \texttt{twin}). Cette feuille de route fixe l'ordre d'implémentation, la répartition des rôles, les jalons datés et la stratégie de plan B.

Trois principes guident le planning :
\begin{itemize}[nosep]
  \item \textbf{Le noyau est testé avant l'affichage.} Une fonction \texttt{checkIntegrity()} vérifie les invariants de la structure (cycles \texttt{next}, symétrie des \texttt{twin}, cohérence des faces et des sommets). Elle est appelée après chaque opération topologique.
  \item \textbf{L'annulation est prévue dès le jalon 3.} Chaque opération d'édition est conçue comme une commande (patron Command) dès son écriture, afin que le jalon 4 soit un branchement et non une refonte.
  \item \textbf{La documentation avance avec le code.} Chaque jalon se termine par sa section du rapport \LaTeX{} et ses captures d'écran, et par des \emph{commits} réguliers, ce qui sert directement les critères de régularité et de documentation.
\end{itemize}

\section{Ordre d'implémentation et jalons}

Les dates ci-dessous correspondent au diagramme de Gantt du projet (GanttProject), démarré le 21 septembre 2026.

\begin{table}[H]
\centering
\small
\begin{tabularx}{\textwidth}{@{}l l Y Y@{}}
\toprule
\textbf{Dates} & \textbf{Jalon} & \textbf{Livrable vérifiable} & \textbf{Critère « terminé »} \\
\midrule
21/09--25/09 & M0 & Dépôt GitHub (Yohan Fougerolle invité), diagramme UML, cette feuille de route, projet CMake qui compile & Un « hello triangle » s'affiche ; la bibliothèque graphique est choisie \\
28/09--15/10 & M1 & Import/export \texttt{.obj} ou \texttt{.ply}, construction du Half-Edge, \texttt{checkIntegrity()}, tests unitaires & Un tétraèdre, un cube et un maillage réel passent import $\to$ export $\to$ import sans perte \\
16/10--26/10 & M2a & Rendu filaire et ombré, caméra orbitale & Un maillage chargé est visible et manipulable \\
28/10--16/11 & M2b & Sélection de sommets, d'arêtes et de faces & Le clic sélectionne le bon élément, mis en évidence \\
17/11--04/12 & M3 & Création de face, découpe d'arête (\emph{edge split}), bascule d'arête (\emph{edge flip}), suppression, le tout sous forme de commandes & Chaque opération passe \texttt{checkIntegrity()} sur des cas intérieurs et de bord \\
07/12--15/12 & M4 & Undo/Redo (Ctrl+Z / Ctrl+Y), séparation données / rendu / interaction, extension originale & Une série de 20 opérations peut être défaite puis refaite ; aucun plantage sur les scénarios de démonstration \\
16/12--24/12 & Rendu & Rapport \LaTeX{}, manuel avec captures, dépôt propre & Un tiers compile et lance le projet en suivant le manuel \\
\bottomrule
\end{tabularx}
\caption{Jalons du projet (dates du diagramme de Gantt, année 2026).}
\end{table}

\section{Rôles et responsabilités}

\begin{table}[H]
\centering
\small
\begin{tabularx}{\textwidth}{@{}l Y@{}}
\toprule
\textbf{Rôle} & \textbf{Responsabilités} \\
\midrule
Noyau (personne A) & Structure Half-Edge, import/export, opérations topologiques, tests d'intégrité \\
Interface (personne B) & Rendu, caméra, sélection, partie interaction de l'undo/redo \\
Partagé & Relecture croisée par \emph{pull requests}, rapport réparti par section, mise à jour du dépôt \\
\bottomrule
\end{tabularx}
\caption{Répartition des rôles (si travail en équipe).}
\end{table}

En travail individuel, ces deux rôles restent des blocs de travail distincts dans le planning, afin de garder une séparation claire entre le noyau et l'interface.

\section{Stratégie de plan B}

Pour chaque risque identifié, une \textbf{date de bascule} est fixée à l'avance, rattachée à la fin du jalon concerné. Si le critère n'est pas atteint à cette date, on passe au plan B sans discussion, afin d'éviter les impasses et les pertes de temps. Ces dates ne figurent volontairement pas dans le diagramme de Gantt principal, pour ne pas le surcharger ; elles sont consignées ici comme registre de suivi séparé.

\begin{table}[H]
\centering
\small
\begin{tabularx}{\textwidth}{@{}Y Y l Y@{}}
\toprule
\textbf{Risque} & \textbf{Plan A} & \textbf{Bascule} & \textbf{Plan B} \\
\midrule
Mise en place de la bibliothèque graphique & OpenGL + GLFW & Fin M0 (25/09) & SFML ou une autre bibliothèque plus simple à installer \\
Reconstruction des \texttt{twin} & Table associative $(v_{\mathrm{origine}}, v_{\mathrm{destination}}) \to$ demi-arête & Fin M1 (15/10) & Refuser les arêtes non-manifold à l'import ; gérer le bord par \texttt{twin = nullptr} \\
Sélection par lancer de rayon & Ray casting 3D & Fin M2b (16/11) & Picking par couleur (rendu hors écran) ou sélection par index dans une liste \\
Cas de bord des opérations (split, flip, suppression) & Gérer bord et intérieur & Milieu M3 (\ensuremath{\sim}26/11) & Traiter d'abord l'intérieur ; l'opération renvoie \texttt{false} si ses préconditions ne sont pas remplies, sans corrompre le maillage \\
Suppression d'un sommet & Réparation locale (remplissage du trou) & Fin M3 (04/12) & Supprimer d'abord toutes les faces incidentes, ou n'autoriser que les sommets isolés \\
Undo/Redo & Commande avec opération inverse & Fin M4 (15/12) & Instantané du maillage avant chaque opération (plus lourd mais fiable) \\
Retard global & Périmètre complet & Fin M3 (04/12) & Réduction du périmètre (section suivante) \\
\bottomrule
\end{tabularx}
\caption{Registre des risques et plans B (dates rattachées aux jalons du Gantt).}
\end{table}

\subsection*{Prototypes courts sur les points incertains}

Avant de s'engager sur un bloc risqué (sélection, suppression de sommet), on réalise un prototype d'une durée maximale de 1 à 2 jours. S'il échoue, le coût est faible et le plan B est déjà identifié.

\subsection*{Périmètre coupable}

Si le temps manque, on sacrifie dans cet ordre :
\begin{enumerate}[nosep]
  \item \textbf{Indispensable} : M1, M2, découpe d'arête, bascule d'arête, rapport.
  \item \textbf{Souhaitable} : création de face, suppression, undo/redo.
  \item \textbf{Bonus} : extensions géométriques, second format de fichier, extension originale.
\end{enumerate}

\end{document}
